% =================================================
% Справочный лист
% =================================================

\begin{center}
  {\fontsize{25}{30}\selectfont\bfseries\sffamily\color{Primary}
    Длина отрезка}

  \vspace{0.3em}

  {\large\sffamily\color{GrayText}
    Измерение, единицы длины и построение}
\end{center}

\vspace{0.45em}

\section{Отрезок и его длина}

\begin{propertybox}[Главная мысль]
  У каждого отрезка есть \textbf{длина}. Длина показывает, во сколько раз
  выбранная единица длины укладывается в данном отрезке.
\end{propertybox}

\segmentlengthfigure

\begin{warningbox}[Не путайте]
  \textbf{Отрезок} — геометрическая фигура, а \textbf{длина отрезка} — число
  с единицей измерения. Запись \(AB=6\text{ см }4\text{ мм}\) сообщает длину
  отрезка \(AB\).
\end{warningbox}

Длины равных отрезков одинаковы. И наоборот: если длины двух отрезков
одинаковы, то эти отрезки равны.

\section{Как измерить отрезок}

\begin{propertybox}[Алгоритм]
  \begin{enumerate}
    \item Совместите один конец отрезка с делением \(0\) линейки.
    \item Расположите линейку вдоль отрезка.
    \item Определите деление, с которым совпадает второй конец.
    \item Запишите число вместе с единицей длины.
  \end{enumerate}
\end{propertybox}

\correctmeasurementfigure

\begin{notebox}[Если отрезок начинается не у нуля]
  Длина равна разности показаний у правого и левого концов.
\end{notebox}

\shiftedmeasurementfigure

В данном примере:
\[
  CD=6\text{ см }8\text{ мм}-1\text{ см }2\text{ мм}
     =5\text{ см }6\text{ мм}.
\]

\begin{warningbox}[Типичная ошибка]
  Нельзя считать край линейки делением \(0\). Начало отсчёта задаёт именно
  штрих с числом \(0\).
\end{warningbox}

\section{Единицы длины}

\unitsladderfigure

\[
  1\text{ км}=1000\text{ м},\qquad
  1\text{ м}=10\text{ дм}=100\text{ см}=1000\text{ мм},
\]
\[
  1\text{ дм}=10\text{ см},\qquad
  1\text{ см}=10\text{ мм}.
\]

\begin{propertybox}[Переход между единицами]
  При переходе к \textbf{меньшим} единицам числовое значение увеличивается;
  при переходе к \textbf{большим} — уменьшается.
\end{propertybox}

Примеры:
\[
  5\text{ см }8\text{ мм}=58\text{ мм},
\]
\[
  905\text{ см}=9\text{ м }5\text{ см},
\]
\[
  6\text{ км }45\text{ м}=6045\text{ м}.
\]

\section{Как построить отрезок заданной длины}

\constructionalgorithmfigure

\begin{propertybox}[Алгоритм построения]
  \begin{enumerate}
    \item Отметьте первый конец отрезка и обозначьте его.
    \item Совместите с ним деление \(0\) линейки.
    \item У нужного деления отметьте второй конец.
    \item Соедините точки и обозначьте второй конец.
  \end{enumerate}
\end{propertybox}

\begin{notebox}[Точность построения]
  Точки должны находиться у нужных штрихов линейки, а линия должна соединять
  их центры. Построение выполняют остро заточенным карандашом.
\end{notebox}

\section{Равные отрезки}

\equalcontrastingsegmentsfigure

На рисунке \(AB=CD\), поэтому отрезки \(AB\) и \(CD\) равны. Отрезок
\(MN\) имеет другую длину.

\section{Точка внутри отрезка}

\partssegmentfigure

\begin{propertybox}[Сложение длин]
  Если точка \(C\) лежит между точками \(A\) и \(B\), то
  \[
    AB=AC+CB.
  \]
\end{propertybox}

Если известны длина всего отрезка и одной его части, вторую часть находят
вычитанием:
\[
  CB=AB-AC.
\]

\partsproblemfigure

\[
  CB=9\text{ см }5\text{ мм}-4\text{ см }8\text{ мм}
    =4\text{ см }7\text{ мм}.
\]

\section{Краткая памятка}

\begin{center}
  \renewcommand{\arraystretch}{1.38}
  \begin{tabularx}{\textwidth}{
    >{\raggedright\arraybackslash}p{0.31\textwidth}
    >{\raggedright\arraybackslash}X}
    \toprule
    \rowcolor{TableHeader}
    Действие & Что нужно сделать \\
    \midrule
    Измерить отрезок
      & совместить один конец с нулём и прочитать деление у второго конца \\
    \addlinespace[0.2em]
    Перевести в меньшие единицы
      & использовать соотношение единиц и выполнить умножение \\
    \addlinespace[0.2em]
    Построить отрезок
      & отметить первый конец, отложить длину и отметить второй конец \\
    \addlinespace[0.2em]
    Найти недостающую часть
      & из длины всего отрезка вычесть длину известной части \\
    \bottomrule
  \end{tabularx}
\end{center}

\begin{questionsbox}[Проверьте себя]
  \begin{enumerate}
    \item Чем отрезок отличается от его длины?
    \item Почему измерение удобно начинать с деления \(0\)?
    \item Как перевести сантиметры в миллиметры?
    \item Как построить отрезок заданной длины?
    \item Как найти одну часть отрезка, если известны весь отрезок и другая часть?
  \end{enumerate}
\end{questionsbox}
